% element: 10-s085:e03
\BeleskeID{10-s085}
Daljim porastom ulaznog napona vodi tranzistor T1, dok je tranzistor T1 zakočen. T1 provodi kompletnu struju $I_E$.

Kolektorski i emiterski napon tada su
\[V_{C1}=V_{CC}-R_{C1}I_E,\qquad V_{E1}=V_I-V_{BE1}=V_I-V_{BE}.\]
Oduzimanjem dobijamo $V_{CE1}=V_{CC}-R_{C1}I_E-V_I+V_{BE}$. Za aktivni režim mora da važi
\[V_{CC}-R_{C1}I_E-V_{CES}+V_{BE}>V_I.\]
Ulazni napon raste, dok se kolektorski napon više ne menja, pa ovaj uslov može prestati da važi. Na prikazanoj granici T1 ulazi u zasićenje; izlazni napon potom prati emiterski napon uvećan za $V_{CES}$. To je neregularna oblast rada ECL kola.
